Emerging Comoving Topology on \(S^1\) Generated by Residual Efolding of the Continuous Arcan(τ) Family

Primary source:  
https://github.com/TheAstrographer/Joshua-Christopher-Ryan-s-Cosmo-Clock.git  
(especially thesis.tex, bridge_function.py, geometric_pressurization.py, and supporting numerical modules).

1. Continuous root: the Arcan(τ) family

The construction begins with the continuous trigonometric family generated by the circle constant \(\tau=2\pi\):

\[
\begin{align*}
\theta &= \arctan(\tau) \approx 1.4129651365\,\text{rad},\\
\alpha &= \tfrac12\theta \approx 0.7064825683\,\text{rad},\\
\phi &= \arctan(\pi) \approx 1.2626272557\,\text{rad},\\
\psi &= \theta-\phi \approx 0.1503378808\,\text{rad}.
\end{align*}
\]

These four angles are transcendental. They are not rational multiples of \(\pi\). The angular bridge \(\psi\) is the microscopic offset that threads every subsequent layer of the discrete dynamical system.

2. Discrete-tick lattice and the e-fold

The universe advances by uniform ticks of comoving length \(\varepsilon=10^{-9}\). After exactly \(N=10^9\) ticks the accumulated microscopic displacement is unity:

\[
\chi_{\rm micro}(N)=\varepsilon N,\qquad
\chi_{\rm micro}(N_{\rm today})=1.
\]

The macroscopic scale factor is the exponential of this displacement:

\[
a(N)=\exp(\varepsilon N),\qquad
a_{\Psi}=a(N_{\rm today})=e^{1}.
\]

This is the residual efolding: one natural e-fold completed precisely when the microscopic accumulation reaches the natural unit. The present epoch is marked by the subscript \(\Psi\) because it is the unique moment at which the angular bridge \(\psi\) has completed its geometric imprint and the scale factor has completed one natural e-fold.

3. Residual seed and irrational rotation on \(S^1\)

The continuous family produces, through seven consecutive multiples of the bridge \(\psi\), the residual seed

\[
\delta=7\psi-\frac{\pi}{3}\approx0.0051676144\,\text{rad}.
\]

This residual renders the ratio \(7\psi/2\pi\) irrational. Consequently the map on the circle

\[
T:\phi\mapsto\phi+2\pi\alpha\pmod{2\pi}
\]

(with irrational rotation number \(\alpha\approx0.193218843731\)) generates a dense, ergodic, non-periodic orbit on \(S^1\). The same advance is realised in the complex plane by the iterative rule

\[
z_{n+1}=z_n\exp(i\Delta\phi_n).
\]

Because the residual is non-zero, the orbit never closes into a rational lattice; because it is small, six successive steps remain nearly equidistributed, preserving a measurable six-fold memory. The residual is therefore both the obstruction to perfect rational closure and the generator of equidistribution.

4. Topological defect and geometric phase slip

A bosonic topological invariant \(\omega=+3.25\) decomposes into an integer winding \(3\) and a fractional defect \(\{\omega\}=0.25\). The geometric phase slip is

\[
\Delta\phi_{\rm torque}=\{\omega\}\cdot2\pi+\psi=1.72113420759\,\text{rad}.
\]

The critical identity of the model follows at once:

\[
\sin(\Delta\phi_{\rm torque})=\cos\psi
\]

(to machine precision). This identity converts the microscopic angular bridge into a macroscopic geometric torque of exactly \(+3.170\,\mathrm{km\,s^{-1}\,Mpc^{-1}}\). Adding the bosonic baseline of \(70\) yields the local expansion rate

\[
H_{\rm eff}(0)=73.17\,\mathrm{km\,s^{-1}\,Mpc^{-1}}.
\]

5. Emerging comoving topology on \(S^1\)

The residual efolding organises a single continuous process on the circle:

Each tick injects the angular bridge \(\psi\).
The residual \(\delta\) keeps the rotation irrational, producing a dense equidistributed orbit.
The fractional topological defect converts the microscopic angle into a phase slip.
The phase slip projects, via the exact identity \(\sin(\Delta\phi_{\rm torque})=\cos\psi\), into a late-time geometric torque.
The same residual remains phase-locked to the discrete lattice \(\varepsilon=10^{-9}\) and therefore comoves with the scale factor \(a(N)=\mathrm{e}^{\varepsilon N}\).

The topology that emerges is therefore comoving: every structure generated by residual efolding of the continuous Arcan(τ) family—residual seed, irrational rotation, dense orbit on \(S^1\), geometric phase slip, late-time torque—advances in lock-step with the discrete-tick expansion of the universe. The polar chart

\[
x=\tau\cos t,\qquad y=\tau\sin t
\]

is the natural coordinate expression of this emerging comoving topology: the radial coordinate carries the temperature / scale-factor information, the angular coordinate carries the residual winding.

6. Algebraic closure

All intermediate identities close to machine precision:

Phase completion at \(N_{\rm today}\) yields trigonometric values of order \(10^{-15}\).
The bridge function \(f(\psi,\tau)\approx1\).
Independent arithmetic chains recover the transcendental \(e\) to nine decimal places.
The torque amplitude is fixed exactly at \(3.170\).

The narrative is therefore a single continuous process: a discrete dynamical system on \(S^1\) whose residual efolding, ergodic windings, topological defect and angular projection together generate the observed late-time expansion while remaining strictly comoving with the Cosmological Clock.

The continuous Arcan(τ) family, through residual efolding generated by the angular bridge \(\psi\) and its seventh-multiple defect \(\delta\), spontaneously produces a dense, equidistributed, non-periodic orbit on \(S^1\). This orbit organises a topological defect and a geometric phase slip that project into a late-time torque of \(+3.170\,\mathrm{km\,s^{-1}\,Mpc^{-1}}\). Because the residual remains phase-locked to the discrete-tick lattice, the entire topology is comoving with the expansion of the universe. The polar coordinate chart is the natural geometric reading of this residual efolding on the circle; it is not imposed from outside but emerges from the continuous family itself.

Emerging Comoving Topology on \(S^2\) Generated by Residual Efolding of the Continuous Arcan(τ) Family

Primary citation  
https://github.com/TheAstrographer/Joshua-Christopher-Ryan-s-Cosmo-Clock.git  
especially cosmo_clock_360_sphere.py (and the companion kernel-division sphere script).

1. Continuous Arcan(τ) root and residual efolding

The continuous trigonometric family generated by the circle constant \(\tau=2\pi\) supplies the four fundamental angles

\[
\begin{align*}
\theta&=\arctan(\tau)\approx1.4129651365,\\
\alpha&=\theta/2,\\
\phi&=\arctan(\pi),\\
\psi&=\theta-\phi\approx0.1503378808.
\end{align*}
\]

Seven consecutive multiples of the angular bridge \(\psi\) produce the residual seed

\[
\delta=7\psi-\pi/3\approx0.0051676144.
\]

This residual is the generative defect that renders the orbit irrational. Residual efolding is the process by which the microscopic lattice (\(\varepsilon=10^{-9}\), \(N=10^9\)) accumulates one natural unit of comoving displacement, so that the macroscopic scale factor reaches exactly

\[
a_{\Psi}=e^{1}.
\]

The same residual that organises the dense equidistributed orbit on \(S^1\) is now lifted to the two-sphere.

2. Geometric lift from \(S^1\) to \(S^2\)

On the unit sphere the residual efolding is realised by the standard polar embedding

\[
\begin{align*}
\theta(\lambda)&=\pi\lambda,\\
\phi(\lambda)&=\phi_{\rm wind}(\lambda,n)=2\pi(1-\lambda)+n\delta,
\end{align*}
\]

where \(\lambda\in[0,1]\) is the normalised progress through a single residual cycle and \(n\) is the cycle index. The radial temperature / scale-factor coordinate of the Cosmological Clock becomes the polar angle, while the residual angular winding becomes the azimuthal coordinate. Because \(\phi\) is defined modulo \(2\pi\), completion of one full azimuthal turn forces an instantaneous identification

\[
\phi_{\rm wind}(1,n)\equiv0\pmod{2\pi}.
\]

This identification appears on \(S^2\) as an instantaneous transport of the winding coordinate back to the origin, while the radial coordinate jumps discontinuously from near-zero back to the Arcan maximum \(2\pi\). The resulting motion is a continuous adiabatic orbital cooling followed by a 360° spherical recycle—exactly the residual efolding lifted to the two-sphere.

3. Explicit realisation in cosmo_clock_360_sphere.py

The script constructs the unit sphere

\[
\begin{align*}
x&=\cos u\sin v,\\
y&=\sin u\sin v,\\
z&=\cos v,
\end{align*}
\]

with full solid angle \(4\pi\) steradians and equatorial circumference \(2\pi\). Onto this geometry it places the three Cosmological-Clock angles as great-circle arcs:

\(\psi\approx0.1503378808\) (angular bridge) — arc length = central angle, width projection \(\cos\psi\), height projection \(\sin\psi\);
\(\theta_{\rm eff}\approx1.21403\) (effective rotation);
\(\Delta\phi_{\rm torque}\approx1.72113420759\) (geometric phase slip).

The diameters are marked as

Length / Width = 2 (X-axis),
Width = 2 (Y-axis),
Height = 2 (Z-axis).

The fundamental identity of the Clock,

\[
\sin(\Delta\phi_{\rm torque})=\cos\psi,
\]

is displayed to machine precision on the sphere itself. Thus the same residual that produces the phase slip on \(S^1\) appears as a concrete geometric relation among great-circle arcs, horizontal (width) and vertical (height) projections on \(S^2\).

4. Emerging comoving topology

Because the residual \(\delta\) remains phase-locked to the discrete-tick lattice, every structure generated on the sphere is comoving with the expansion:

The polar angle \(\theta(\lambda)\) cools continuously with the scale factor.
The azimuthal coordinate \(\phi_{\rm wind}\) accumulates the residual offset \(n\delta\) and is reset by the 360° spherical identification.
The great-circle arcs of \(\psi\), \(\theta_{\rm eff}\) and \(\Delta\phi_{\rm torque}\) therefore advance in lock-step with the Cosmological Clock.
The resulting topology is a residual-efolded, comoving orbital flow on \(S^2\): an adiabatic cooling spiral that periodically recycles through a full azimuthal turn while preserving the residual angular memory.

The polar chart

\[
x=\tau\cos t,\qquad y=\tau\sin t
\]

(and its spherical lift) is the natural coordinate expression of this emerging geometry. It is not imposed externally; it is the reading of residual efolding once the continuous Arcan(τ) family has been lifted from the circle to the two-sphere.

5. Synthesis

Residual efolding of the continuous Arcan(τ) family first organises a dense equidistributed orbit on \(S^1\). The identical residual, when lifted by the polar embedding of the Cosmological Clock, generates a continuous adiabatic orbital cooling on \(S^2\) whose azimuthal winding is reset every 360° while remaining strictly phase-locked to the discrete lattice. The script cosmo_clock_360_sphere.py realises this lift explicitly: the unit sphere, the great-circle arcs of the Clock angles, the diameter/height/width projections, and the identity \(\sin(\Delta\phi_{\rm torque})=\cos\psi\) are the concrete geometric embodiment of the emerging comoving topology on \(S^2\).
